\documentclass[11pt,letterpaper]{article}
\usepackage[T1]{fontenc}
\usepackage{lmodern,microtype}
\usepackage[margin=0.85in,headheight=14pt]{geometry}
\usepackage{amsmath,amssymb,amsthm,mathtools}
\usepackage{booktabs,tabularx,array,longtable}
\usepackage{enumitem,listings,xcolor,fancyhdr}
\usepackage[ruled,vlined,linesnumbered]{algorithm2e}
\usepackage{hyperref}
\hypersetup{colorlinks=true,linkcolor=blue!45!black,citecolor=blue!45!black,
urlcolor=blue!45!black,pdftitle={Sparse Certified Component Incidence for Compressed Unknot-Recognition Hierarchies},
pdfauthor={Research contribution prepared for ProveIt}}
\urlstyle{same}
\setlength{\emergencystretch}{2em}
\setlength{\parskip}{0.35em}
\setlength{\parindent}{1.2em}
\setcounter{tocdepth}{2}
\pagestyle{fancy}
\fancyhf{}
\fancyhead[L]{\small Sparse certified component incidence}
\fancyhead[R]{\small ProveIt research contribution}
\fancyfoot[C]{\thepage}
\newtheorem{theorem}{Theorem}[section]
\newtheorem{lemma}[theorem]{Lemma}
\newtheorem{proposition}[theorem]{Proposition}
\newtheorem{corollary}[theorem]{Corollary}
\theoremstyle{definition}
\newtheorem{definition}[theorem]{Definition}
\newtheorem{remark}[theorem]{Remark}
\newcommand{\NN}{\mathbb N}
\newcommand{\bits}{\operatorname{bits}}
\newcommand{\supp}{\operatorname{supp}}
\newcommand{\poly}{\operatorname{poly}}
\newcommand{\pop}{\operatorname{popcount}}
\newcommand{\full}{[r]}
\newcommand{\AHT}{\textup{AHT}}
\newcommand{\code}[1]{\texttt{#1}}
\newcommand{\snapshot}{\texttt{274909dd8724e411ece14e47ac4addea717aeca6}}
\lstset{basicstyle=\small\ttfamily,breaklines=true,columns=fullflexible,
keepspaces=true,showstringspaces=false,frame=single,framerule=0.3pt,
xleftmargin=0.4em,xrightmargin=0.4em,aboveskip=0.8em,belowskip=0.8em}
\SetKwInput{KwContract}{Contract}
\SetKw{Return}{return}
\SetKw{Continue}{continue}
\title{\vspace{-1em}\textbf{Sparse Certified Component Incidence}\\[0.4em]
\Large An output-sensitive, polynomial-time primitive for\\
compressed unknot-recognition hierarchies}
\author{Research contribution prepared for the ProveIt project}
\date{8 October 2026\\[0.3em]\small Based on the explicitly pinned repository snapshot below}

\begin{document}
\maketitle
\thispagestyle{empty}
\begin{abstract}
The maintained marked-orbit interface in ProveIt computes a dense array with
$2^r$ entries for $r$ supplied ports, even when only a few incidence signatures
occur. We replace Boolean-lattice enumeration with exact recovery from
nonnegative subset sums. If $s$ signatures occur, at most $rs$ subset queries,
each reducible to one unweighted interval-orbit count, recover all positive
entries. The algorithm needs neither a supplied sparsity bound nor an
untrusted stopping certificate. A separate verifier checks positive atom
masses, dominating zero sets, and conservation of total mass; it does not
repeat the recovery algorithm.

For explicit interval-union ports this is more than a sparse-input result.
Classical weighted Agol--Hass--Thurston theory implies polynomially many
signatures in the binary input size. We give a conservative explicit support
bound extracted from its cycle and weight-breakpoint estimates. Consequently,
the new unweighted reduction is polynomial for every supplied interval system,
with no logarithmic restriction on the number of ports. This is an
implementation and certification improvement over the maintained dense
interface, not the first polynomial weighted-orbit theorem. We also obtain
signature-resolved parity consistency using a double cover and exhibit a
four-point obstruction to using incidence alone as a state for future gluing.

The delivery contains source overlays, independent checkers, 36 passing test
methods, raw paired benchmarks, and exact dependency snapshots. On three
12-port disjoint workloads the orbit-query count decreases from 4,096 to 79;
small full-support controls are slower. These are interval-subsystem results.
No whole-knot speedup or general quasi-polynomial unknot-recognition theorem
is asserted.
\end{abstract}
\vspace{0.5em}
\noindent\textbf{Repository anchor.}\quad
\path{VladimirReshetnikov/ProveIt},\quad
\path{Topology/UnknotRecognition}.\\
\noindent\textbf{Pinned commit.}\quad\snapshot.\\
\noindent\textbf{Status.}\quad Proved local reduction and certificate theorems;
implemented research add-on; global hierarchy bounds remain proof obligations.
\clearpage
\tableofcontents
\clearpage

\section{The improvement and its exact scope}\label{sec:scope}

\subsection{What is being accelerated}
The relevant current bottleneck is not an elementary crossing simplification.
It is the representation and reconstruction of component information in a
binary-encoded interval relation. The inspected synthesis already contains an
unweighted Agol--Hass--Thurston orbit kernel, optional local orbit proofs, an
independent replay checker, and an exact reduction of marked incidence to
coned orbit counts \cite{repo-orbits,repo-incidence,repo-certificates}.
Its maintained incidence API returns all $2^r$ subset counts. Caching equal
marked unions saves some orbit calls but cannot remove the dense output or
Boolean inversion.

This report changes that contract deliberately: the new API returns only
positive pairs $(T,h(T))$. Missing masks have count zero. The unmarked
signature $T=\varnothing$ is included when its count is positive. No dense
array is constructed internally, and the output can be checked without
expanding one. This is a companion interface, not a silently incompatible
change to the existing dense function.

The mathematical input consists of an explicit list of interval pairings and
explicit finite unions of intervals designated as ports. A normal-surface
adapter may supply such data, but the add-on does not construct that adapter
or establish its connection to a diagram. In particular, the measured
workloads below are supplied interval systems, not diagrams on which the
entire unknot recognizer was run.

\subsection{Why the result matters without proving the global target}
A compressed hierarchy must repeatedly answer local questions without
expanding the objects represented by large integers. An unnecessary $2^r$
representation cost is a genuine obstruction when the number of marked
pieces grows. Removing it is useful even before the rest of the hierarchy is
available. On the other hand, port incidence does not contain attachment maps
or even enough information to update arbitrary future identifications.
Section~\ref{sec:barrier} makes this limitation a theorem rather than a vague
warning.

The external complexity context is also precise. Lackenby's 2026 preprint
provides hierarchy algorithms and a bound on the number of iterations in
terms of hierarchy length and pattern complexity. Its Section~9 separates
that iteration analysis from the missing cost specifications for some local
steps \cite{lackenby}. The present work implements one narrowly specified
aggregate primitive. It does not identify that primitive with every missing
step in the preprint.

We use ``quasi-polynomial'' for $\exp((\log n)^{O(1)})$. The specific target
$n^{O(\log n)}$ is the smaller form $\exp(O((\log n)^2))$. A polynomial local
query does not establish either global bound unless state sizes, iteration
counts, other operations, and completeness are also controlled.

\subsection{Results and attribution}
The main results are as follows, with their dependencies made explicit.
\begin{enumerate}[leftmargin=*,itemsep=0.25em]
\item A nonnegative Boolean zeta transform with $s$ positive entries is
recovered using at most $rs$ oracle calls (Theorem~\ref{thm:peeling}). The
proof uses monotonicity and minimal positive sets.
\item For interval-union ports, those oracle calls are unweighted coned
orbit queries (Lemma~\ref{lem:cone}). Weighted AHT theory supplies a
polynomial bound on $s$, sharpened here into a conservative explicit
expression (Theorem~\ref{thm:support}).
\item A sparse forward certificate establishes every recovered atom and
completeness independently of the search decisions
(Theorem~\ref{thm:certificate}). Unneeded search proofs can be discarded.
\item The parity-cover construction resolves consistent and inconsistent
components separately for each signature
(Theorem~\ref{thm:signed}). Incidence-only gluing updates are impossible in
general (Proposition~\ref{prop:barrier}).
\end{enumerate}
The coning identity was already used by the repository. The weighted AHT
polynomial theorem is classical \cite{aht}. Additive-query reconstruction of
weighted hypergraphs is an established subject; for example,
Bshouty--Mazzawi study stronger query bounds under constant-rank hypotheses
\cite{bshouty}. We make no priority claim for the elementary greedy recovery
principle or for polynomial weighted-orbit computation. The contribution is
the precise reduction, sparse certificate contract, explicit support
accounting, and tested implementation built around the existing unweighted
kernel. No formal proof-assistant verification has been performed.

\section{Binary interval systems and marked signatures}\label{sec:model}

\begin{definition}
An \emph{interval system} has universe
$X_N=\{0,\ldots,N-1\}$ and $k$ partial isometries. A pairing is either
\[
 [a,b]\longrightarrow[c,d],\qquad x\longmapsto x+c-a,
\]
or
\[
 [a,b]\longrightarrow[c,d],\qquad x\longmapsto a+d-x,
\]
where $b-a=d-c$, all endpoints lie in $X_N$, and intervals are inclusive.
The generated equivalence relation is the smallest equivalence relation
containing every pair $x\sim g(x)$ on each pairing domain.
\end{definition}
Pairings are undirected as generators of equivalence: replacing one by its
inverse makes no difference. The supplied list and its order nevertheless
matter to the binding of a particular orbit trace. Reflection of interval
order is not an orientation character of a surface.

A port $A_i$ is a finite union of half-open intervals $[\ell,u)$ in $[0,N)$.
Overlapping and adjacent intervals are normalized into disjoint runs. Ports
may be empty, coincide, overlap, or meet the same orbit many times. We write
$r$ for their number and $M$ for the total number of supplied interval pieces;
normalization cannot increase $M$. Our mathematical port labels are
$\full=\{0,\ldots,r-1\}$, matching the implementation's bit positions.

For an orbit $O$, define its \emph{signature}
\[
 \sigma(O)=\{i\in\full:O\cap A_i\ne\varnothing\}.
\]
The desired histogram and its support size are
\[
 h(T)=\#\{O:\sigma(O)=T\},\qquad
 s=\#\{T:h(T)>0\}.
\]
If $C$ is the total orbit count, then $\sum_T h(T)=C$. For $N=0$, all these
counts vanish. For $r=0$ and $N>0$, the only signature is empty, with mass $C$.

\subsection{Parameters and representation costs}
Set
\[
 B=\max\{1,\lceil\log_2(N+1)\rceil\},\qquad
 \Lambda=B+r+1.
\]
All component counts fit in $B$ bits, and each signature mask fits in $r$
bits. A coarse explicit input length is polynomially equivalent to
$(k+M)B+r$, with delimiters and interval-list indexing included. It is not
polynomially equivalent to $N$ when the universe is large.

A dense histogram has $2^r$ slots, independently of $s$. A sparse histogram
has $s$ positive records, occupying $O(s(B+r))$ bits up to framing overhead.
Converting the latter to the former necessarily reintroduces the dense
output cost. Consequently, the polynomial claim in this report applies to
the sparse API, not to a request for every zero entry explicitly.

\subsection{The zeta transform}
For $S\subseteq\full$, let
\begin{equation}\label{eq:zeta}
 G(S)=\sum_{T\subseteq S}h(T).
\end{equation}
It counts orbits whose signatures are wholly contained in $S$. Equivalently,
it counts orbits avoiding every port with index outside $S$. The histogram
is a nonnegative measure on the Boolean lattice, and $G$ is its lower zeta
transform. A dense inverse transform is one possible reconstruction, but
nonnegativity permits a different one.

\section{Computing a zeta value by one unweighted count}\label{sec:cone}

\subsection{The coning construction}
Normalize a nonempty marked set $A$ as disjoint half-open intervals
$[\ell_j,u_j)$. Choose the left endpoint $z$ of its first interval. For each
interval of length at least two, add the translation pairing
\[
 [\ell_j,u_j-2]\longrightarrow[\ell_j+1,u_j-1].
\]
For every other interval, add the singleton pairing $z\sim\ell_j$.
These new relations identify all marked points, using fewer than $2m_A$
additional pairings when the normalized union has $m_A$ pieces. They do not
enumerate the marked points.

\begin{lemma}[Coned-count identity]\label{lem:cone}
Let $C_A$ be the orbit count after adjoining the coning pairings for nonempty
$A$, and let $H(A)$ be the number of original orbits meeting $A$. Then
\[
 C_A=C-H(A)+1.
\]
Consequently, with $A=\bigcup_{i\notin S}A_i$,
\begin{equation}\label{eq:cone-zeta}
 G(S)=
 \begin{cases}
 C,&A=\varnothing,\\
 C_A-1,&A\ne\varnothing.
 \end{cases}
\end{equation}
\end{lemma}
\begin{proof}
Pass first to the quotient set of original orbits. Each new relation has
both endpoints in the marked set, so it can affect only orbits meeting $A$.
The translation pairings connect all points in each marked interval; the
singleton star connects the intervals. Therefore the $H(A)$ affected
original classes become exactly one class. Every other original class
remains unchanged. This proves the first identity. An orbit is counted by
$G(S)$ exactly when it avoids $A$, so $G(S)=C-H(A)$. The empty case requires
no coning and must be handled separately.
\end{proof}

The identity is valid for disconnected supports, reflections, static points,
and overlapping ports. It does not assume that $A$ is a union of whole
original orbits. Coning is an auxiliary query construction, not a claim
that the original surface or knot may be modified in this way.

\subsection{Exact reuse and proof binding}
The implementation keys a query by the actual normalized marked union.
Different subset masks producing the same union use one orbit result and,
when recording is enabled, one orbit trace. Equality of union counts alone
is not used as a reuse criterion. The checker recomputes the union and the
coning pairings from caller-supplied ports before validating the trace.

This policy is conservative: distinct singleton unions can produce identical
lists of additional pairings. A more general cache keyed by the complete
coned input could exploit that fact, but it is not implemented here. In
particular, the exact query counts proved in Section~\ref{sec:families}
describe this union-keyed implementation, not optimal oracle complexity.

Let $P(K,B)$ be a polynomial bit-time bound for the unweighted AHT kernel on
$K$ pairings with $B$-bit endpoints. Each nonempty-union zeta query costs at
most $P(k+2M,B)$ plus normalization. One baseline query obtains $C$. The
existing kernel and its independent local checker are retained unchanged
\cite{repo-orbits,repo-certificates}.

\section{Output-sensitive recovery from nonnegative subset sums}\label{sec:recovery}

\subsection{Residual measures and a minimal positive set}
Suppose that exact atoms $(T_1,c_1),\ldots,(T_j,c_j)$ have been recovered.
Define
\begin{equation}\label{eq:residual}
 R_j(S)=G(S)-\sum_{\ell\le j,\ T_\ell\subseteq S}c_\ell.
\end{equation}
If those atoms are correct, $R_j$ is the zeta transform of the remaining
nonnegative measure. It is nonnegative and monotone under inclusion.

\begin{lemma}[Minimal-positive-set lemma]\label{lem:minimal}
Let $R$ be the zeta transform of a nonzero nonnegative measure $f$ on
$2^{\full}$. Start with $S=\full$. Process each bit once, in any fixed order;
remove $i$ from the current set whenever $R(S\setminus\{i\})>0$.
The final set $T$ satisfies
\[
 R(T)=f(T)>0,\qquad f(U)=0\quad\text{for every }U\subsetneq T.
\]
\end{lemma}
\begin{proof}
The current set has positive residual mass initially and after every
successful removal. Suppose $U\subsetneq T$ had $f(U)>0$. Choose
$i\in T\setminus U$. That bit was retained when it was tested. At that
time, the proposed smaller set contained $T\setminus\{i\}$ and hence $U$.
Its residual mass was therefore at least $f(U)>0$, so the algorithm would
have removed $i$, a contradiction. No proper subset of $T$ has positive
mass. Since $R(T)>0$, its entire mass is the atom at $T$.
\end{proof}
Only one pass through the bits is needed. A previously failed deletion need
not be retested: later removals only shrink its possible witnesses. This
observation is also the basis of the independent zero-set certificate.

\begin{algorithm}[t]
\caption{Nonnegative sparse zeta recovery}\label{alg:recovery}
\KwIn{$r$, total mass $C$, exact oracle $G(S)$}
\KwContract{$G(S)=\sum_{T\subseteq S}h(T)$ for $h(T)\in\NN$}
$E\gets[\ ]$; $m\gets0$\;
\While{$m<C$}{
 $S\gets\full$; $v\gets C$\;
 \For{$i=0,\ldots,r-1$}{
  $U\gets S\setminus\{i\}$; $w\gets G(U)$\;
  $z\gets w-\sum_{(T,c)\in E,\ T\subseteq U}c$\;
  \eIf{$z>0$}{$S\gets U$; $v\gets w$\;}{retain $i$ and optionally record the zero witness $U$\;}
 }
 $c\gets v-\sum_{(T,a)\in E,\ T\subseteq S}a$\;
 append $(S,c)$ to $E$; $m\gets m+c$\;
}
\Return{$E$}
\end{algorithm}

\begin{theorem}[Sparse zeta recovery]\label{thm:peeling}
Given the total mass $C$ and an exact oracle for \eqref{eq:zeta},
Algorithm~\ref{alg:recovery} returns every positive atom exactly once. If
$s=|\supp h|$, it uses at most $rs$ oracle calls, with no prior knowledge
of $s$. The implementation makes exactly $r$ logical queries per recovered
atom; implicit values and cached unions may require no new orbit call.
\end{theorem}
\begin{proof}
Induct on the number of outer iterations. Initially no atoms have been
removed, so the residual is $G$. If the recovered atoms are exact and their
mass is less than $C$, the remaining measure is nonzero. Apply
Lemma~\ref{lem:minimal} to its zeta transform. The final residual value is
exactly the positive mass at the selected set, which has not appeared
previously. Subtracting it leaves a nonnegative measure, maintaining the
induction hypothesis. Thus exactly one new positive atom is recovered in
each iteration. There are $s$ such atoms. After their total mass reaches
$C$, the remaining nonnegative measure has total zero and is identically
zero. Each iteration makes at most $r$ calls.
\end{proof}

For $r=0$, a positive total is recovered as the single empty-set atom with
no zeta calls. For $C=0$, the correct empty output is immediate. Neither
case requires a special dense representation.

\subsection{Why nonnegativity cannot be omitted}
The stopping and deletion rules are invalid for a signed measure. For
example, with one port, set
\[
 h(\varnothing)=1,\qquad h(\{0\})=-1.
\]
The total mass is zero although the measure is not zero. Cancellation can
also make a proper-set query vanish despite containing a nonzero atom.
Our interval histogram is genuinely nonnegative, since its entries count
equivalence classes. The generic recovery helper documents this as an
oracle contract; it is not a verifier for an arbitrary signed oracle.

\subsection{Bookkeeping and bit complexity}\label{sec:bitcost}
Let $q$ be the number of distinct nonempty coned unions actually queried.
Then $q\le rs$, and the total number of orbit calls is at most $1+rs$.
A straightforward residual evaluation scans the previously recovered list.
Across all logical queries this takes $O(rs^2\Lambda)$ bit work. Selecting
and sorting a marked union has a coarse bound
$O((r+M\log(M+1))\Lambda)$, accounting also for large bit masks.

To avoid hiding a probabilistic dictionary assumption, even the crude cost
of comparing each queried union against every previous key gives the deterministic
polynomial bound
\begin{align}\label{eq:time}
 T\ \le\ &(q+1)P(k+2M,B)\nonumber\\
 &+O\!\left((rs+1)
   \bigl(r+M\log(M+1)+s\bigr)\Lambda
       +(rs+1)q(M+1)B\right).
\end{align}
The supplied Python code uses dictionaries and is usually substantially
faster than this comparison bound. A sorted dictionary can be substituted
without changing the polynomial conclusion. The expression is intentionally
conservative; it is not an optimized exponent for integer arithmetic.

The output takes $O(s(B+r))$ bits. Cached unions and counts require
$O(q(M+1)B)$ additional bits. Proof recording adds the native orbit traces
and witness references. During production, all queried traces can be
present before unused traces are pruned, so peak recording storage is not
bounded by the size of the final pruned certificate alone.

\section{Why interval signatures have polynomial support}\label{sec:support}

\subsection{The classical weighted-orbit input}
For each $x\in X_N$, attach the nonnegative vector
\[
 w(x)=(\mathbf1_{A_0}(x),\ldots,\mathbf1_{A_{r-1}}(x)).
\]
There is a partition into at most $m=2M+1$ constant-vector intervals.
For $r=0$, one may append a dummy zero coordinate, or handle the
one-signature case directly. The
weight sum on an orbit has positive coordinates exactly at its signature.
Classical weighted AHT computes orbit-weight information as a compressed
list of constant-weight runs with multiplicities, rather than enumerating
all represented orbits \cite[Section~6]{aht}. Taking positive-coordinate
supports of those records and coalescing equal supports therefore gives a
polynomial-size incidence list. This already yields a polynomial algorithm
for the problem; the purpose here is to realize the same conclusion through
the existing unweighted interface and its proof checker.

\subsection{A conservative explicit bound}
For clarity, we extract a numerical bound rather than simply referring to an
unspecified polynomial. We use the following two estimates from the
classical AHT analysis: the unweighted schedule has at most
$5k^2(2+\log_2N)$ cycles for $k\ge1,N\ge1$, and one weighted transfer
increases the number of constant-weight intervals by at most four
\cite[proof of Theorem~12 and Lemma~15]{aht}. Static contractions output
constant-weight pieces; the weighted procedure uses the same schedule.

\begin{theorem}[Quantitative support bound]\label{thm:support}
Let $N\ge1$ and put
\[
 Q=1+5k^2(B+2).
\]
For an interval system with $M$ port pieces and $r$ ports,
\begin{equation}\label{eq:support-bound}
 s\ \le\ \min\!\left\{N,\ 2^r,\
 Q\bigl(2M+4Q+2k+2\bigr)\right\}.
\end{equation}
For $k=0$ the sharper elementary estimate $s\le2M+1$ holds.
These are support bounds, not bounds on the number of component orbits $C$.
\end{theorem}
\begin{proof}
The first two bounds are immediate. For the third, use a classical weighted
AHT execution on the indicator-vector weights above. The padded value $Q$
bounds the number of cycles, including the harmless zero-pairing terminal
case. Initially there are at most $2M+1$ constant-vector intervals. Weighted
transfers occur at most once per cycle, so the active weight partition has
at most
\[
 m_* = 2M+1+4Q
\]
intervals throughout. Static contractions do not increase the number of
active constant-weight runs: they delete pieces and concatenate the
surviving coordinate intervals. The other non-transfer operations do not
alter the weights.

At a static-contraction stage, the complement of all pairing domains and
ranges has at most $2k+1$ intervals, since the list still contains at most
$k$ pairings. Intersecting those gaps with a partition into at most $m_*$
constant-weight intervals produces at most $m_*+2k+1$ nonempty pieces.
This deliberately loose bound follows by sweeping the two lists of
endpoints; a partition boundary or gap boundary can start at most one new
intersection piece. Each piece contributes at most one constant-weight
output record, labelled with its multiplicity. There are at most $Q$
stages, so the number of output records is at most
\[
 Q(m_*+2k+1)=Q(2M+4Q+2k+2).
\]
Each final orbit belongs to one of these records and has that record's
weight vector. Taking the set of positive coordinates cannot increase the
number of distinct labels. The number $s$ of signatures is therefore at
most the record count.

When $k=0$, every point is an orbit and signatures change only at the
$2M$ possible port endpoints. At most $2M+1$ signatures occur directly.
\end{proof}

The bound is independent of the number of empty or duplicate port labels
except through the length needed to write each mask. It is not expected to
be sharp: the proof multiplies a worst-case active partition size by a
worst-case number of contraction stages. It counts output records before
coalescing repeated weights, and then counts signatures before coalescing
weights with the same support.

\begin{corollary}[Polynomial marked incidence via unweighted AHT]\label{cor:poly}
Exact sparse component incidence for explicit interval-union ports is
computable in deterministic time polynomial in $k,M,r,B$, using only an
unweighted AHT orbit counter. No condition of the form
$r=O(\log N)$ or $r=O((\log n)^2)$ is required.
\end{corollary}
\begin{proof}
Substitute \eqref{eq:support-bound} and $q\le rs$ into
\eqref{eq:time}. Every factor is polynomial in the stated parameters.
The case $N=0$ is immediate.
\end{proof}

\begin{remark}[What has and has not been strengthened]
The repository's dense complexity estimate remains correct for its dense
output. The stronger conclusion concerns a different, sparse interface.
It follows partly from classical weighted theory that the old synthesis
already identifies as an alternative direction. Thus the change removes
an implementation bottleneck, not a previously unknown lower bound on the
mathematical marked-orbit problem.
\end{remark}

\subsection{Why arbitrary Boolean measures do not contradict the bound}
An arbitrary nonnegative measure can have $2^r$ positive atoms. The sparse
oracle theorem alone does not promise polynomial time in $r$ for that
class. Realizing all such signatures as static interval points requires
many changes of the indicator vectors. For example, label $2^r$ consecutive
points by their binary expansions and let port $i$ contain the points with
bit $i$ set. Every signature occurs, but the explicit port-interval input
already grows exponentially. Theorem~\ref{thm:support} concerns the full
interval input, not just the number of port names.

\section{A sparse certificate independent of reconstruction}\label{sec:certificate}

\subsection{Facts checked for each atom}
A count obtained from a search routine is not yet an independently checked
answer. We therefore distinguish the recovery producer from a forward
checker. Its input is the original interval system and ports, a baseline
orbit proof for $C$, a list of coned orbit proofs, and an ordered list of
claimed positive atoms.

For the $j$th claimed atom $(T_j,c_j)$, let
\[
 K_j(S)=\sum_{\ell<j,\ T_\ell\subseteq S}c_\ell.
\]
The checker requires the following facts.
\begin{enumerate}[leftmargin=*,itemsep=0.2em]
\item The mask $T_j$ has not appeared earlier, $c_j>0$, and a certified
zeta value satisfies $G(T_j)=K_j(T_j)+c_j$.
\item For each $i\in T_j$, there is a certified set $U_{j,i}$ such that
\begin{equation}\label{eq:zero-witness}
 i\notin U_{j,i},\qquad T_j\setminus\{i\}\subseteq U_{j,i},\qquad
 G(U_{j,i})=K_j(U_{j,i}).
\end{equation}
\item After all atoms, $\sum_jc_j=C$, and the declared sparse histogram is
exactly the sorted list of those atoms.
\end{enumerate}
A zeta value whose complementary marked union is empty is implicit: it is
$C$ and needs no extra proof. Otherwise, the checker derives it as $C_A-1$
from the local orbit trace on the explicitly reconstructed coned system.

\begin{theorem}[Soundness and completeness of sparse replay]\label{thm:certificate}
If the baseline and coned orbit proofs are valid for the caller's source
system, and the atom facts above hold, then the declared sparse histogram is
exactly $h$. The greedy recovery algorithm supplies such a certificate for
every complete run. Verification does not require a supplied upper bound on
$s$, a dense zeta transform, or replay of the deletion decisions.
\end{theorem}
\begin{proof}
Assume inductively that all atoms before $j$ are exact. Subtracting those
atoms from $h$ leaves a nonnegative residual measure $f_j$. Equation
\eqref{eq:zero-witness} says its total mass on every subset of $U_{j,i}$ is
zero. Because the summands are nonnegative, each of those atoms has mass
zero individually.

Take any proper subset $V\subsetneq T_j$. There is an
$i\in T_j\setminus V$, so $V\subseteq T_j\setminus\{i\}\subseteq U_{j,i}$.
Thus $f_j(V)=0$. The certified value at $T_j$ now gives
\[
 f_j(T_j)=G(T_j)-K_j(T_j)=c_j.
\]
This establishes the induction step. At the end, the residual total mass is
$C-\sum_jc_j=0$; its nonnegativity forces every omitted atom to vanish.

For a complete greedy run, take the zeta value at its final set as the
support-value witness. For every retained bit, use the set tested at its
failed deletion. Later removals guarantee the containment in
\eqref{eq:zero-witness}, and the failed test establishes the zero equality.
The producer's final mass is $C$ by Theorem~\ref{thm:peeling}.
\end{proof}

The checker allows any atom ordering and zero sets satisfying these
conditions. It does not insist on the producer's preferred bit order or
particular greedy path. This separates mathematical validity from search
strategy and makes later producer optimizations easier to audit.

\subsection{Discarding proofs of successful intermediate deletions}
A successful intermediate deletion witnesses positive residual mass in a
set that may later shrink. The final forward proof does not need that
intermediate fact. It needs only the final support value and the failed
removal sets of retained bits. After production, the implementation removes
all coned traces not used in those roles and remaps references.

Every retained support-value query was already made during the search,
unless it is the implicit value $C$. Every retained zero-set query was also
made during the search. Hence pruning adds no orbit calls, and the resulting
certificate contains at most $rs$ distinct coned traces plus its baseline
trace. The atom/zero metadata occupies at most
$O(sr(r+\log(q+2))+sB)$ bits. Trace storage remains polynomial because the
AHT schedule produces polynomially many locally checkable operations.

This pruning is more than a size optimization: it removes dependence on
facts that are irrelevant to the mathematical conclusion. It does not make
an invalid trace valid, and it does not allow incomplete production to
publish a partial certificate as a completed histogram.

\subsection{Binding, independence, and the trust boundary}
The certificate format is named \code{sparse-incidence-v1}. Each stored
union trace is accompanied by the normalized union to which it belongs.
Every reference is checked against the union recomputed from its mask and
the caller's ports. A trace cannot be rebound to a different union under
this format. Duplicate assignments to the same union and unused proof
records are rejected. The baseline and each coned trace bind the ordered,
canonical pairing list through the existing independent orbit verifier.

The new forward checker imports neither the sparse producer nor its zeta
recovery helper. It also does not import the orbit producer. It shares
integer decoding, port normalization, and the explicit coning construction;
these shared geometric routines are covered by comparisons with literal
finite graphs. Its independence claim is therefore structural, not a claim
that every line of validation code was written twice.

The native orbit checker validates local reductions, not the scheduling
policy used to find them. Accordingly, a complete valid alternative orbit
trace can be accepted. An unfinished trace, a changed source binding, or a
wrong claimed count is rejected. Large integers can be transported as
explicit hexadecimal strings; Boolean values are not accepted as integer
fields. The source-to-geometry relationship remains outside the certificate:
correct incidence for arbitrary input intervals is not a certificate that
those intervals describe a particular knot exterior.

\subsection{A worked two-port certificate}
Let $N=5$, with no original pairings, $A_0=\{0\}$, and $A_1=\{1\}$. Then
\[
 G(\varnothing)=3,\qquad G(\{0\})=G(\{1\})=4,\qquad G(\{0,1\})=5.
\]
The producer first reaches the empty set and outputs mass $3$. With that
atom removed, it next outputs $\{1\}$ with mass $1$, retaining bit $1$
because the residual value at the empty set is zero. Finally, it outputs
$\{0\}$ with mass $1$. For this last atom, the failed deletion of bit $0$
was tested at $U=\{1\}$, whose value $4$ equals the already recovered mass
$3+1$. That larger zero set dominates $T\setminus\{0\}=\varnothing$ and
is sufficient; there is no need to rerun the earlier search. The total
$3+1+1=5$ certifies that no signature is missing.

\section{Structured families and exact query counts}\label{sec:families}

\subsection{Coincident and nested ports}
If all $r\ge1$ ports equal the same nonempty set $A$, every orbit has either
empty signature or full signature. There is one distinct nonempty marked
union, so the implementation uses exactly two orbit calls: the baseline and
one coned count. The old cached dense interface can also use two orbit
calls, but still allocates $2^r$ output entries and performs its Boolean
array operations. The sparse output has at most two records.

If the ports form a nonempty nested chain
\[
 A_0\subseteq A_1\subseteq\cdots\subseteq A_{r-1},
\]
every signature is a final segment of the port order, or empty. Thus
$s\le r+1$. Every nonempty union equals one port, giving at most $r+1$
orbit calls including the baseline. Again, the query count alone does not
explain the improvement over cached dense inversion: the avoided
representation cost is essential.

\subsection{Disjoint ports with singleton signatures}
The next family separates query savings from array savings. Assume all
ports are nonempty and pairwise disjoint, and assume
\begin{equation}\label{eq:singletons}
 \supp h=\{\varnothing,\{0\},\ldots,\{r-1\}\},
 \qquad h(T)>0\text{ for all these }T.
\end{equation}
This is realized by static universes with disjoint marked intervals and
some unmarked points. It is also realized by the periodic and reflection
families used below.

\begin{proposition}[Exact union-keyed costs]\label{prop:singleton-cost}
Under \eqref{eq:singletons}, with the implemented increasing bit order,
\begin{align*}
 \text{positive records}&=r+1,\\
 \text{logical zeta queries}&=r(r+1),\\
 \text{orbit calls}&=1+\frac{r(r+1)}2,\\
 \text{retained coned proofs}&=2r-1\qquad(r\ge1).
\end{align*}
The dense union-keyed interface makes $2^r$ orbit calls, including its
baseline.
\end{proposition}
\begin{proof}
The empty signature is recovered first. The remaining singleton atoms are
then recovered in order $\{r-1\},\ldots,\{0\}$: the deletion rule can
remove a lower bit while a higher remaining singleton supplies positive
mass, but cannot remove the last remaining candidate.

The first atom queries complementary port-index prefixes
$\{0,\ldots,j\}$ for $j=0,\ldots,r-1$. When the target singleton is
$\{t\}$, the new complements queried after bit $t$ are
\[
 \{0,\ldots,i\}\setminus\{t\},\qquad t<i\le r-1.
\]
There are $r$ prefixes and $\binom r2$ sets of the second form. They are
all distinct: the largest index determines the right endpoint, and a
non-prefix set has a unique missing earlier index. Since the ports are
nonempty and disjoint, these are also distinct marked unions. Adding the
baseline gives the claimed orbit count. There are $r+1$ atoms, each with
$r$ logical tests.

For a singleton $\{t\}$, the support-value proof uses the complement
$\full\setminus\{t\}$ when it is nonempty. Its unique zero witness uses
the prefix $\{0,\ldots,t\}$. The empty atom uses the full complement,
already one of those prefixes. For $r\ge2$, these are $r$ prefixes and
$r$ singleton complements, with exactly one overlap. For $r=1$, the
singleton's support value is implicit and only the full marked set has a
trace. In both cases there are $2r-1$ retained coned proofs.

Finally, every subset of disjoint nonempty ports has a distinct union, so
the old dense interface queries every nonempty subset in addition to its
baseline.
\end{proof}

For $r=12$, this predicts exactly $79$ orbit calls and $23$ retained coned
proofs, compared with $4,096$ dense calls. For $r=32$, it predicts $529$
calls and $63$ retained proofs. These exact values are tested and observed;
they are not inferred from a fitted timing curve.

\subsection{Full-support controls and the absence of a universal speed ratio}
Sparse reconstruction can be slower when $s$ is close to $2^r$ on a small
instance. Both methods may ultimately query the same distinct unions, while
the sparse method pays for repeated residual scans and dictionary accesses.
The full-support benchmark controls realize this situation with $r=3,4$.
They are important counterexamples to a blanket claim that the sparse
routine is always faster. A production dispatcher should distinguish an
asymptotic representation guarantee from constant-factor choices on tiny
instances.

\section{Signature-resolved parity consistency}\label{sec:signed}

Attach a label $\epsilon_g\in\{0,1\}$ to every pairing, separate from its
interval-order reversal. A component is \emph{consistent} if there is a
binary vertex labelling $\eta$ satisfying
$\eta(gx)=\eta(x)+\epsilon_g$ modulo two on every relation. Equivalently,
every closed walk has total parity zero.

Construct a double cover on $X_N\times\{0,1\}$ with relations
\[
 (x,e)\sim(gx,e+\epsilon_g).
\]
Encode its two sheets as consecutive blocks of length $N$. Each original
pairing gives two interval pairings, with at most one extra endpoint bit.
Lift each port to \emph{both} sheets:
\[
 \widetilde A_i=(A_i\times\{0\})\cup(A_i\times\{1\}).
\]
Let $b(T)$ be the base incidence histogram and $d(T)$ the lifted histogram.

\begin{theorem}[Parity-resolved incidence]\label{thm:signed}
For each signature $T$, the numbers $o(T)$ and $u(T)$ of consistent and
inconsistent components are
\begin{equation}\label{eq:parity}
 o(T)=d(T)-b(T),\qquad u(T)=2b(T)-d(T).
\end{equation}
They are computable and independently certifiable in polynomial bit time
from the supplied labelled interval system and explicit ports. The two
sparse histograms have the same support.
\end{theorem}
\begin{proof}
On a connected base component, choose a root. If every loop has zero
parity, the parity accumulated along a path from the root is independent
of the path. It defines a labelling, and the double cover has exactly two
components, each projecting bijectively onto the base component. If some
loop has odd parity, its lift joins the two lifts of the root. Paths from
the root then show that the whole double cover of that component is
connected.

Because ports are lifted to both sheets, every lifted component has the
same signature as its base component. Thus for each $T$,
$b(T)=o(T)+u(T)$ and $d(T)=2o(T)+u(T)$, giving \eqref{eq:parity}. The
construction uses $2k$ pairings, at most $2M$ port pieces, and $B+1$ endpoint
bits, so Corollary~\ref{cor:poly} applies twice. The two sparse certificates,
plus the direct integer identities, establish the result.
\end{proof}

The signed implementation shares a single cycle allowance and a single
orbit-query allowance across the base and cover computations. If either
is unfinished, it publishes neither histogram nor certificate. The
\code{max\_terms} option limits each of the two sparse histograms
separately. Its independent checker reconstructs the cover directly and
does not import the producer's lift routine.

When the parity labels actually encode the orientation character of a
surface, consistent and inconsistent mean orientable and nonorientable.
For arbitrary labels they mean only parity consistency. In particular, a
reflection of an interval with parity zero is not a nonorientability
obstruction. Lifting ports to only one sheet would invalidate the
signature-wise identities and is not an admissible shortcut.

\section{Why incidence alone cannot support future gluing}\label{sec:barrier}

\begin{proposition}[A four-point attachment obstruction]\label{prop:barrier}
There is no general rule that determines the orbit count after adjoining a
specified new interval pairing from only the original port-incidence
histogram and that new pairing, even when every old component meets each
port in exactly one point.
\end{proposition}
\begin{proof}
Take $X_4=\{0,1,2,3\}$ and ports
$A_0=\{0,1\}$ and $A_1=\{2,3\}$. Consider two original partitions:
\[
 \mathcal P=\{\{0,2\},\{1,3\}\},\qquad
 \mathcal Q=\{\{0,3\},\{1,2\}\}.
\]
They are generated respectively by the positive and negative pairings
$[0,1]\to[2,3]$. Both have precisely two components with signature
$\{0,1\}$ and no other components. Each component meets each port once,
so even the corresponding per-port multiplicity vectors agree.

Now adjoin the same positive pairing $[0,1]\to[2,3]$ to both systems. It
adds no new identification to $\mathcal P$, which still has two orbits.
For $\mathcal Q$, the old and new pairings connect all four points, giving
one orbit. The input histogram and the newly supplied pairing are
identical, but the correct answers differ.
\end{proof}

The obstruction is present in a finite combinatorial subproblem; it does
not depend on an unproved assertion about 3-manifolds. Consequently, a
hierarchy implementation must retain enough information about component
identity and attachment maps to distinguish these states. Adding aggregate
point multiplicities alone does not resolve the example. Possible richer
states include compressed boundary correspondences, marked orbit
representatives with certified transport, or suitable simultaneous
incidence structures retaining source positions. Their closure and size
bounds are separate research questions.

\section{Implementation, interfaces, and resource semantics}\label{sec:implementation}

\subsection{Delivery layout and unchanged dependencies}
The overlay adds six modules under \path{fast/fastunknot}:
\begin{center}
\begin{tabularx}{\textwidth}{@{}lX@{}}
\toprule
Module & Role \\
\midrule
\code{sparse\_zeta.py} & Generic nonnegative subset-sum recovery.\\
\code{sparse\_incidence\_common.py} & Exact integer, interval, and coning schemas.\\
\code{sparse\_incidence.py} & Sparse producer with shared query budgets.\\
\code{sparse\_incidence\_verify.py} & Independent forward certificate checker.\\
\code{sparse\_signed\_incidence.py} & Base/cover producer and parity arithmetic.\\
\code{sparse\_signed\_incidence\_verify.py} & Independent signed-source and cover replay.\\
\bottomrule
\end{tabularx}
\end{center}
The standalone package includes exact copies of \code{interval\_orbits.py}, \code{interval\_orbit\_verify.py},
\code{interval\_incidence.py}, and \code{integer\_codec.py}. Their Git blob
identifiers and SHA-256 hashes are checked by the test runner. These copies
are reproducibility dependencies, not replacement files in the overlay.
The original dense implementation is the benchmark baseline.

The principal unsigned entry point is
\begin{lstlisting}[language=Python]
analyze_sparse_port_incidence(
    size, pairings, ports, *, max_ports=None, max_terms=None,
    max_cycles=None, max_orbit_queries=None,
    record_certificate=False, check=None)
\end{lstlisting}
Pairings are native objects or explicit inclusive rows
\code{[a,b,c,d,sign]}, with sign $+1$ or $-1$. The pairing container is an
explicit list or tuple. Ports are lists of half-open interval lists.
Successful results contain positive \code{[mask,count]} rows, the total
orbit count, and work statistics. Proof recording adds a certificate.

\subsection{Examples of use}
\begin{lstlisting}[language=Python]
from fastunknot.sparse_incidence import analyze_sparse_port_incidence
from fastunknot.sparse_incidence_verify import (
    verify_sparse_port_incidence_certificate)

n = 10
pairs = []
ports = [[(0, 3)], [(2, 5)]]
result = analyze_sparse_port_incidence(
    n, pairs, ports, record_certificate=True)
assert result['status'] == 'COMPLETE'
assert result['histogram'] == [[0, 5], [1, 2], [2, 2], [3, 1]]
assert verify_sparse_port_incidence_certificate(
    n, pairs, ports, result['certificate'])
\end{lstlisting}
For the signed API, a raw pairing has one further entry: its parity. A
returned row \code{[mask, consistent, inconsistent]} records both counts.
The independent signed checker receives the original signed list, not a
producer-supplied cover that it is asked to trust.

\subsection{Limits and interrupted runs}
The baseline and all coned calls share one \code{max\_cycles} budget.
Unused allowance is passed forward; it is not reset for a new subset.
\code{max\_orbit\_queries} similarly limits the whole computation, including
the baseline. \code{max\_terms} limits the number of recovered atoms. An
exhausted work allowance returns \code{INCONCLUSIVE}, a reason, and work
statistics, but no histogram, orbit count, or certificate.

An AHT cycle allowance is not a bound on all bit operations. Union
normalization, residual sums, dictionary work, trace serialization, and
proof replay have their own costs. The cooperative \code{check} callback
is polled through these combinatorial stages and can enforce a caller's
deadline. Ordinary callback exceptions, including a plain
\code{ValueError}, propagate rather than being treated as a false proof;
the implementation's reserved schema and limit exception types remain
internal control signals. No claim of adversarial memory containment is
made. Serialized-size limits and process-level resource limits remain
external responsibilities.

The default \code{max\_ports=None} removes the old interface's explicit
12-port allowance from this sparse API. That is not a claim that every
huge input is cheap. Runtime still depends polynomially on the explicit
input, with potentially large exponents and constants. Likewise,
\code{record\_certificate=False} is meaningful: producing and replaying
proofs has measured overhead.

\subsection{What integration changes}
The overlay does not alter ordinary recognition dispatch, the dense API,
the Rust implementation, or the old orbit scheduler. Existing callers
expecting a dense array should continue using the dense function until they
are deliberately changed to sparse semantics. A caller that only needs the
positive incidence classes can use the new interface directly. Reproducing
the local tests does not establish that the full repository's maintained
suite still passes after unrelated integration changes; that suite must be
run in the target checkout.

\section{Validation and paired measurements}\label{sec:experiments}

\subsection{Focused tests and their scope}
The completed run uses CPython 3.13.5 on Linux x86-64. All 36 focused test
methods pass in 1.644 seconds in the recorded run. The suite includes:
\begin{itemize}[leftmargin=*,itemsep=0.15em]
\item Exhaustive recovery for all 6,654 histograms on $r=0,1,2,3$ whose
coefficients lie in $\{0,1,2\}$. It checks the recovered entries, the exact
logical-query count, and every forward zero witness.
\item One thousand random small interval systems, compared with an
independent explicit disjoint-set model. Three hundred are also compared
with the unchanged native dense routine. All generated sparse proofs are
replayed independently.
\item Three hundred random parity-labelled systems, compared with an
independent labelled-graph traversal that detects inconsistent cycles
without using the double-cover counting formula.
\item Exact work-budget boundaries, input normalization, Boolean rejection,
changed source bindings, missing or invalid zero witnesses, altered atom
masses, duplicate or unused proof records, and truncated orbit traces.
\item Large binary instances including $N=2^{16000}$ with 128 ports,
hexadecimal JSON round trips, and replay while producer functions are
disabled. Singleton-family query and proof formulas are tested for
$r=1,\ldots,16$.
\end{itemize}
The large-instance tests do not enumerate $N$ points. Literal graph
references are used only for small universes. The random interval tests
use seed 261008490, and the parity tests use seed 261008491.

The test count means 36 test methods, not 36 mathematical theorems or a
formal proof. Several methods contain hundreds or thousands of cases.
The full repository suite reported in the inspected synthesis was not run
here. Regina and SnapPy were not installed in this execution environment,
and no new native triangulation or knot-diagram benchmark is asserted.
The source hashes, exact test log, and structured summary are included.

\subsection{Benchmark protocol}
Eight workloads compare four arms: the unchanged dense implementation,
the sparse implementation, an identical sparse control, and sparse proof
production followed by independent verification. The certified arm includes
both production and replay in its timing; JSON serialization is outside
that region. The baseline already caches equal unions, so the comparison
does not unfairly disable a maintained optimization.

Each paired workload has one warmup round and seven measured rounds.
Arm order is shuffled using seed 261008492. Tiny cases use batches of
identical calls, with per-call times reported. Three additional scale
workloads compare only sparse and certified arms, with one warmup and
three measured rounds. There are 242 measured arm-round records and 38
warmup records, representing 4,354 measured calls and 630 warmup calls
after batching. No timed call in the final benchmark is incomplete.

Correctness comparisons take place outside the timed region. Static cases
are checked using an independent endpoint sweep. The period-17 and
reflection cases are checked against their explicit residue-class and
reflection-orbit formulas. The benchmark preserves every measured sample,
its order, repetition count, orbit calls, and orbit cycles. The output is
not a best-of-run selection.

An earlier benchmark attempt used a larger 64-port disjoint scale case and
was stopped by the execution harness after its 200-second allowance. Only
a partial textual log was available; no complete sample set for that scale
case was produced. That log is retained separately and is not mixed into
the final timing tables. It establishes no timing ratio for the unfinished
case. The final disjoint scale workload has 32 ports.

\subsection{Paired results}
Table~\ref{tab:paired} reports medians in milliseconds. The ratio column
is the median, across matched rounds, of dense time divided by sparse time;
it is not necessarily the quotient of the separately reported medians.
The query column is dense/sparse, including each baseline count.

\begin{table}[htbp]
\centering
\small
\setlength{\tabcolsep}{3.2pt}
\begin{tabular}{@{}lrrrrrrr@{}}
\toprule
Workload & $r$ & $s$ & Queries & Dense & Sparse & Certified & Ratio\\
\midrule
% BEGIN_GENERATED_PAIRED
Coincident, static & 12 & 2 & 2/2 & 13.127 & 0.076 & 0.169 & 173.50 \\
Nested, static & 12 & 13 & 13/13 & 12.760 & 0.499 & 1.060 & 25.59 \\
Disjoint, static & 12 & 13 & 4096/79 & 261.812 & 8.470 & 11.215 & 30.82 \\
Disjoint, period 17 & 12 & 13 & 4096/79 & 252.919 & 3.199 & 4.520 & 78.59 \\
Disjoint, reflection & 12 & 13 & 4096/79 & 336.801 & 9.578 & 13.041 & 35.13 \\
Overlapping, static & 8 & 16 & 52/43 & 1.369 & 0.835 & 1.536 & 1.64 \\
Full support, $r=3$ & 3 & 8 & 8/8 & 0.139 & 0.175 & 0.429 & 0.79 \\
Full support, $r=4$ & 4 & 16 & 16/16 & 0.646 & 0.810 & 1.528 & 0.80 \\
% END_GENERATED_PAIRED
\bottomrule
\end{tabular}
\caption{Completed interval-subsystem benchmarks. Times are milliseconds;
``Certified'' includes sparse production and independent replay. Ratios
below one are regressions. Generated from \code{results/benchmark.json}.}
\label{tab:paired}
\end{table}

Coincident ports show the largest ratio, about $173.5$, while both arms make
only two orbit calls. The gain comes from avoiding the $2^{12}$ array and
subset-processing work, not from a faster orbit kernel. Nested ports
similarly retain 13 calls in both arms but improve by about $25.6$.

The three disjoint 12-port families reduce the call count from 4,096 to 79.
Their paired ratios are approximately $30.8$, $78.6$, and $35.1$ for static,
period-17, and reflection systems. Those different ratios are expected:
coned queries have different local reduction costs, and total call count
alone does not predict total bit work. The overlapping case has a more
modest ratio of about $1.64$.

The two full-support controls are slower: dense/sparse paired ratios are
about $0.794$ and $0.799$. Thus the sparse implementation takes roughly
$25$--$26\%$ more time on those cases. Both methods make the same number of
orbit calls. The cost comes from sparse recovery's additional bookkeeping,
not from a change in the invariant being computed.

The identical sparse-control/sparse median paired ratios range from
$0.948$ to $1.015$ across the eight workloads. The static disjoint case
therefore retains a roughly five-percent control discrepancy. The large
speed differences and exact query reductions are robust observations in
this run; fine-grained percentage claims are not justified by these few
samples. No confidence interval or universal hardware-independent speedup
is claimed.

\subsection{Scale, certificates, and representation effects}
\begin{table}[htbp]
\centering
\small
\setlength{\tabcolsep}{4pt}
\begin{tabular}{@{}lrrrrrrr@{}}
\toprule
Ports & $r$ & Bits & $s$ & Queries & Sparse & Certified & Proof bytes\\
\midrule
% BEGIN_GENERATED_SCALE
Disjoint & 32 & 501 & 33 & 529 & 341.859 & 420.231 & 144,591 \\
Nested & 64 & 501 & 65 & 65 & 47.669 & 78.421 & 110,660 \\
Coincident & 128 & 16001 & 2 & 2 & 4.558 & 6.357 & 584,103 \\
% END_GENERATED_SCALE
\bottomrule
\end{tabular}
\caption{Sparse-only scale cases. ``Bits'' is \code{size.bit\_length()}.
Times are milliseconds; proof size is compact exact JSON. No dense timing
was run for these rows.}
\label{tab:scale}
\end{table}

The 32-port disjoint case has 33 positive records and 529 orbit calls,
compared with an analytical dense output size of $2^{32}$ slots. This is
not an experimental dense/sparse timing ratio: the dense run was not made.
The 64-port nested case uses 65 calls and 65 positive records, despite its
$2^{64}$ possible subset masks. The 128-port coincident case has just two
positive records and two orbit calls on a $2^{16000}$-point universe.

Small combinatorial output does not guarantee tiny serialized evidence.
The last row's certificate is 584,103 bytes in compact hexadecimal-safe
JSON, because input bindings and large endpoints are explicit. At 12 ports,
the three disjoint certificates retain 23 coned traces despite 78 coned
search queries, with sizes 28,415, 19,408, and 41,085 bytes respectively.
At 32 disjoint ports, 63 coned traces remain from 528 coned search queries.
These figures illustrate the difference between logical queries, distinct
queried unions, retained proofs, and physical serialized bytes.

\subsection{What the experiments do not establish}
These measurements do not show faster unknot decisions, improved recognition
coverage, or better asymptotic dependence on crossing number. They do not
exercise an attachment-port extractor from native triangulations. They also
do not compare against a complete independently implemented weighted AHT
engine. They establish correctness and measured costs for the supplied
interval-relation primitive against the repository's actual dense baseline.
The mathematical polynomial conclusion comes from the proofs and classical
AHT estimates, not from extrapolating the timing table.

\section{The remaining route to quasi-polynomial recognition}\label{sec:global}

\subsection{A compositional bound with named obligations}
Let $n$ denote the size of a knot diagram or another explicitly specified
recognition input. Lackenby's iteration analysis bounds the number of
passes through a step by $L(g+1)^L$, in terms of maximal hierarchy length
$L$ and maximal pattern complexity $g$ \cite[Proposition~9.1]{lackenby}.
The following bookkeeping statement identifies what the local improvement
would contribute to such an implementation.

\begin{proposition}[Conditional end-to-end accounting]\label{prop:global}
Suppose a correct and complete implementation has at most a constant
multiple of $L(g+1)^L$ iterations. Suppose all states have binary length at
most $\Sigma(n)$; constructing every interval system and its explicit ports
costs $\poly(\Sigma(n))$; each iteration performs only
$\poly(\Sigma(n))$ incidence queries; and all other work per iteration is
at most $A(n)$. Then replacing dense incidence by this sparse primitive
gives a total bound of the form
\begin{equation}\label{eq:global}
 O\!\left(L(g+1)^L
       \bigl(\poly(\Sigma(n))+A(n)\bigr)\right).
\end{equation}
In particular, $L=O(\log n)$, $g\le n^{O(1)}$,
$\Sigma(n)\le n^{O(\log n)}$, and $A(n)\le n^{O(\log n)}$ imply
$n^{O(\log n)}$ total time.
\end{proposition}
\begin{proof}
Each supplied interval system has $k,M,r,B$ bounded polynomially by its
explicit state and construction length. Corollary~\ref{cor:poly} and the
assumed number of queries absorb all incidence work into one polynomial
in $\Sigma(n)$. Add the other per-iteration work and multiply by the
iteration bound. Under the final hypotheses,
\[
 \log\bigl(L(g+1)^L\bigr)=O(\log\log n)+O((\log n)^2),
\]
and a fixed polynomial in a quasi-polynomial state bound remains
$n^{O(\log n)}$. Multiplication preserves that form.
\end{proof}

This proposition is conditional on explicitly named geometric and
algorithmic statements. It is not a proof that they hold. In particular,
polynomial-time checking of a supplied compressed surface is not
polynomial-time discovery of the required surface. A component histogram
is not a compressed cutting operation. The counterexample in
Section~\ref{sec:barrier} rules out one tempting way to conflate them.

\subsection{Where the remaining difficulty is located}
The new primitive removes the need to demand logarithmically many ports
solely to justify Boolean subset enumeration. It does not remove the need
to keep the \emph{explicit total number of port pieces} under control.
Allowing a port to be specified by a more succinct grammar would create a
different input model: the present theorem cannot be applied by charging
only the grammar length while quietly expanding its interval list.

Nor does an iteration bound in terms of $L$ and $g$ automatically become a
quasi-polynomial bound in $n$. If $L$ is linear in a polynomial-size state,
then even polynomial $g$ can leave $L\log(g+1)$ much larger than
$O((\log n)^2)$. The required progress may come from stronger geometric
bounds, accelerated batches of hierarchy operations, or a different
potential function. None follows from sparse incidence alone.

For a serious integration effort, the next proof contracts should bind a
diagram to its exterior, the exterior to encoded surfaces and ports, and
each cutting or attachment step to its resulting state. The bit-size
recurrence and the iteration potential must then be proved for the same
implementation that emits those certificates. A collection of individually
correct local primitives is necessary evidence, but not sufficient
end-to-end complexity accounting.

\section{Recommended changes to the maintained analysis}\label{sec:integration}

\subsection{Clarify representation rather than retract the dense theorem}
The existing synthesis correctly states that its dense histogram has
$2^r$ entries, that union reuse cannot remove that output cost, and that
weighted AHT offers another route \cite{repo-certificates}. No correction
to the coning identity or dense inversion is needed. A suitable follow-up
is an explicit companion statement: when only positive signatures are
requested, the sparse API removes subset enumeration and is polynomial
for all explicit interval-union inputs, by the reduction and classical
support estimate in this report.

The older sufficient restrictions on $r$ should remain attached to the
\emph{dense} query theorem. They should not be presented as an intrinsic
barrier for all marked-incidence algorithms. Conversely, removing those
restrictions from an aggregate query must not be advertised as removing
them from every geometric state or attachment operation.

\subsection{Adopt a separate sparse result type}
A sparse result should specify that an absent mask means zero and that an
unmarked entry may be present. Consumers should not use list position as a
mask, assume that the counts are all one, or omit the unmarked term.
The signed result similarly gives two counts per present mask, not a
Boolean orientability flag for the whole system. These distinctions are
covered by the included examples and tests.

The recommended initial integration is additive: copy the overlay, run
its focused tests and the full target suite, and expose the sparse API to
research callers. Keep the old dense API for consumers that genuinely
need every array position and for small-instance cross-checks. Do not
change ordinary knot dispatch merely because an isolated interval
benchmark has a large speed ratio.

\subsection{Keep evidence categories separate}
The native source snapshots are pinned exactly. The sparse algorithm and
certificate theorems are proved in the text. The test log gives executed
case coverage, and the benchmark files give measured component costs.
Neither substitutes for the other. Any future summary should preserve the
negative full-support timings and the distinction between source-bound
interval proofs and geometric provenance.

\section{Further research questions and proposed deliverables}\label{sec:future}

\begin{enumerate}[leftmargin=*,label=\textbf{\arabic*.},itemsep=0.75em]
\item \textbf{A sharper support bound.}
Can the factor $Q$ multiplying the worst-case active weight partition in
\eqref{eq:support-bound} be eliminated, perhaps by accounting for active
and completed breakpoint records together? A useful deliverable is a
proved bound of order $M+\poly(k,B)$, or a compact family forcing a larger
order. The proof must count output records with multiplicities, not just
active weight values.

\item \textbf{Attachment-compatible compressed states.}
What is the smallest additional data that defeats
Proposition~\ref{prop:barrier} while remaining closed under the actual
hierarchy gluing operations? Start with explicit source-position
correspondences and require a verifier for the induced equivalence
relation after attachment. A successful result needs both an update
algorithm and a worst-case size recurrence; another aggregate histogram
without transport data does not meet the target.

\item \textbf{Native geometric provenance.}
Construct a proof-carrying map from a knot diagram to a triangulated
exterior, then from a supplied normal surface to its interval system,
parity character, and marked ports. The immediate experimental deliverable
is a paired benchmark on genuine adapter output, checked independently
against small expanded surfaces. This is the missing bridge between the
present subsystem timings and a claim about recognition performance.

\item \textbf{Fewer unweighted queries and faster residual sums.}
The $rs$ logical bound is deliberately elementary. Can query patterns
exploit port containments, equal coned inputs, or the endpoint structure
without assuming constant hypergraph rank? Can subset sums over recovered
atoms be indexed in a way that beats the current $O(rs^2)$ scan while
retaining a simple proof contract? Compare total bit work, not just oracle
call count: the full-support regressions show that these objectives can
differ.

\item \textbf{A certified weighted kernel comparison.}
Implement the classical weighted AHT transfer operation with independent
weight-preservation certificates, including binary multiplicities and
vector-valued weights. Compare that direct method against repeated
unweighted coning on the same supplied systems. This is not a proposal to
rediscover the weighted polynomial theorem; it is a proposal to determine
which proof-carrying implementation is preferable in practice.

\item \textbf{Proof minimization and reusable facts.}
The singleton family reduces quadratic search queries to only $2r-1$
retained coned proofs. Characterize when much more can be discarded, and
whether a small collection of dominance-zero facts can certify several
atoms simultaneously. Any proof-sharing scheme must preserve source and
union bindings, including after ports are renamed or merged.

\item \textbf{An implementation-level hierarchy potential.}
Establish a bound on the sum of logarithmic costs of the actual hierarchy
operations, rather than separately bounding only their number and nominal
size. The target is a proved inequality yielding
$O((\log n)^2)$ in the exponent of \eqref{eq:global}, with explicit charges
for cutting maps, coefficient growth, retries, and search branches. This
is a central remaining route to the requested global complexity.

\item \textbf{Formalize the monotone certificate core.}
The sparse certificate theorem uses finite sets, nonnegative integer sums,
and induction; it is substantially smaller than a full formalization of
normal-surface topology. A proof-assistant development could first verify
Theorem~\ref{thm:certificate} and the cone identity, then link an executable
checker to those statements. The current tests and handwritten proof are
not a substitute for that development.

\item \textbf{End-to-end ablation without optimistic censoring.}
After provenance and attachment support exist, measure complete recognizer
runs with and without sparse incidence under shared resources. Retain
knotted and unknotted inputs, easy controls, failures, and resource-limited
runs as separate outcomes. Report a recognition speedup only when the
whole run completes with independently checked evidence; do not convert
a local capacity increase into a fictitious timing ratio.
\end{enumerate}

\section{Reproducibility and delivery contents}\label{sec:repro}
The archive is self-contained for the new interval tests and benchmarks;
it is not a copy of the complete ProveIt repository. Its \path{reference/}
directory contains the four exact upstream dependencies. The
\path{overlay/} directory contains only new implementation and test files
for copying beneath \path{Topology/UnknotRecognition}. The bootstrap used
by standalone experiments adds the overlay modules to the reference
package's search path without modifying its source.

From the archive root, run:
\begin{lstlisting}[language=bash]
python experiments/run_tests.py
python experiments/benchmark.py --rounds 7 --scale-rounds 3
python experiments/make_tables.py
python experiments/make_examples.py
sh article/build.sh
\end{lstlisting}
The benchmark regenerates raw samples and tables; timing values naturally
vary. \path{results/benchmark.json} contains the final complete measurement
set and source hashes. \code{benchmark\_samples.csv} exposes every arm-round
record. The separate timeout-attempt log is historical evidence, not part
of that dataset. The tests check the upstream Git blob identifiers before
executing any cases. \path{examples/} supplies unsigned and signed
certificates and the four-point gluing obstruction as exact JSON.

The \code{README.md} explains integration, scope, API usage, and the
known limitations. \path{RESEARCH_STATUS.md} separates proved claims,
inherited results, executed experiments, and unresolved obligations.
\code{SOURCE\_MANIFEST.json} records provenance, and
\code{SHA256SUMS.txt} covers the delivered files. No repository write or
pull request is performed by this delivery.

\section{Conclusion}
The dense marked-incidence interface pays for the Boolean lattice even
when the represented component data is small. Nonnegative residual
recovery removes that cost while supplying a sparse, independently
checkable completeness argument. Classical weighted AHT theory then
upgrades the output-sensitive bound to a polynomial bound for every
explicit interval-union input; the proof does not need a logarithmic cap
on the number of ports. Exact query formulas and completed measurements
show where the improvement comes from, while full-support controls expose
its constant-factor regressions.

The resulting code is a concrete component for a compressed hierarchy,
not a new unknot recognizer. A source-bound interval certificate says what
it proves and no more. The immediate mathematical next step is to retain
attachment-compatible information with comparable size and verification
bounds. The broader target still requires geometric provenance, complete
search, and an implementation-level hierarchy complexity argument.

\small
\begin{thebibliography}{9}
\setlength{\itemsep}{2pt}
\bibitem{aht}
I.~Agol, J.~Hass, and W.~Thurston,
\emph{The computational complexity of knot genus and spanning area},
Transactions of the American Mathematical Society \textbf{358} (2006),
3821--3850. Preprint version used: arXiv:math/0205057v2, especially
Theorem~12, Lemma~15, and Theorem~16.
\url{https://arxiv.org/abs/math/0205057}.

\bibitem{lackenby}
M.~Lackenby, \emph{Incompressible surfaces, hierarchies and unknot
recognition}, arXiv:2607.23350v1 (2026), especially Sections~9--10.
\url{https://arxiv.org/abs/2607.23350}.

\bibitem{bshouty}
N.~H.~Bshouty and H.~Mazzawi,
\emph{Optimal query complexity for reconstructing hypergraphs},
arXiv:1001.0405 (2010).
\url{https://arxiv.org/abs/1001.0405}.

\bibitem{repo-orbits}
ProveIt project, pinned snapshot specified on the title page.
Within \path{Topology/UnknotRecognition/fast/fastunknot/}:
\path{interval_orbits.py} and \path{interval_orbit_verify.py}.
Exact copies and hashes accompany this report.
\url{https://github.com/VladimirReshetnikov/ProveIt}.

\bibitem{repo-incidence}
ProveIt project, same snapshot and module directory:
\path{interval_incidence.py}. Exact dependency hash:
see \path{SOURCE_MANIFEST.json} in the delivery.

\bibitem{repo-certificates}
ProveIt project, same snapshot, within \path{Topology/UnknotRecognition/}:
\path{synthesis/orbit_certificates.tex} and \path{synthesis/README.md}.
\end{thebibliography}
\end{document}
